% c04 原书 p.30–44（PDF p.037–051）
\subsection{概形与态射的性质}
\label{sec:1.3}

\begin{proposition}\label{it:1.3.1}
设 $A$ 是环。以下命题等价：
\begin{enumerate}[label=(\roman*),leftmargin=2.8em]
  \item $\Spec A$ 连通。
  \item $A$ 除 $0$ 与 $1$ 外没有幂等元。
  \item 对每个分解 $A\simeq A_1\times A_2$，都有 $A_1=0$ 或 $A_2=0$。
\end{enumerate}
\end{proposition}

\begin{proof}
(i)$\Rightarrow$(ii) 设 $e$ 是 $A$ 的幂等元。有 $e(1-e)=0$。故
\[
V((e))\cup V((1-e))=V(e(1-e))=\Spec A,
\]
\[
V((e))\cap V((1-e))=V((e+1-e))=\emptyset.
\]
此外 $V((e))$ 与 $V((1-e))$ 是闭集。若 $\Spec A$ 连通，则 $V((e))=\Spec A=V((0))$ 或 $V((1-e))=\Spec A=V((0))$。由 \ref{it:1.2.2}(iv)，$e$ 或 $1-e$ 幂零。但 $e$ 与 $1-e$ 都是幂等元。于是 $e=0$ 或 $1-e=0$。

(ii)$\Rightarrow$(iii) 设 $A=A_1\times A_2$。则 $(1,0)$ 幂等。故 $(1,0)=(0,0)$ 或 $(1,0)=(1,1)$，从而 $A_1=0$ 或 $A_2=0$。

(iii)$\Rightarrow$(i) 假设 $\Spec A$ 不连通。则可找到 $\Spec A$ 的两个非空开子集 $U$ 与 $V$，使
\[
U\cup V=\Spec A,\quad U\cap V=\emptyset.
\]
考虑由 $e|_U=1$ 与 $e|_V=0$ 定义的截面 $e\in\cO_{\Spec A}(\Spec A)$。则 $e$ 幂等。通过同构 $\cO_{\Spec A}(\Spec A)\cong A$，$e$ 定义了 $A$ 中一个异于 $0$ 与 $1$ 的幂等元。令 $A_1=(e)$，$A_2=(1-e)$。则 $A_1$（resp.\ $A_2$）是以 $e$（resp.\ $1-e$）为单位元的环。可以验证 $A\cong A_1\times A_2$。
\end{proof}

非空拓扑空间 $X$ 称为\emph{不可约}\idx{bukeyue}{不可约}的，若它满足以下等价条件之一：
\begin{enumerate}[label=(\roman*),leftmargin=2.8em]
  \item $X$ 不是两个真闭子集的并，即若 $X=F_1\cup F_2$，其中 $F_1$ 与 $F_2$ 是闭子集，则 $F_1=X$ 或 $F_2=X$。
  \item $X$ 的任意两个非空开子集的交非空。
  \item 每个非空开子集都稠密。
\end{enumerate}

注意不可约拓扑空间必连通。设 $X$ 是拓扑空间，$Y$ 是 $X$ 的子集。若 $Y$ 不可约，则 $\ol{Y}$ 不可约。

\begin{proposition}\label{it:1.3.2}
$\Spec A$ 的闭子集不可约，当且仅当它具有形式 $V(\fp)$，其中 $\fp$ 是 $A$ 的素理想。
\end{proposition}

\begin{proof}
设 $\fp$ 是 $\Spec A$ 的素理想。若 $V(\fp)=V(\fa)\cup V(\fb)$，则 $V(\fp)=V(\fa\fb)$，从而 $\fa\fb\subset \sqrt{\fp}=\fp$。由此得 $\fa\subset \fp$ 或 $\fb\subset \fp$，于是 $V(\fp)=V(\fa)$ 或 $V(\fp)=V(\fb)$。这证明了 $V(\fp)$ 不可约。

反之，设 $Y$ 是 $\Spec A$ 的不可约闭子集。可找到 $A$ 的理想 $\fp$，使 $Y=V(\fp)$。由 \ref{it:1.2.2}(ii) 与 $\sqrt{\sqrt{\fp}}=\sqrt{\fp}$，可取 $\fp$ 使 $\sqrt{\fp}=\fp$。设 $\fa\fb\subset \fp$。则
\[
V(\fp)\subset V(\fa\fb)=V(\fa)\cup V(\fb).
\]
因 $V(\fp)$ 不可约，有 $V(\fp)\subset V(\fa)$ 或 $V(\fp)\subset V(\fb)$，从而 $\fa\subset \fp$ 或 $\fb\subset \fp$。所以 $\fp$ 是素理想。
\end{proof}

\begin{proposition}\label{it:1.3.3}
设 $X$ 是概形。对 $X$ 的任意不可约闭子集 $Y$，存在唯一点 $y\in Y$，使 $Y=\ol{\{y\}}$。
\end{proposition}

我们称 $y$ 为 $Y$ 的\emph{一般点}\idx{yibandian}{一般点}。

\begin{proof}
设 $U=\Spec A$ 是 $X$ 的仿射开子概形，使 $U\cap Y$ 非空。则 $U\cap Y$ 是 $U$ 的不可约闭子集，故 $U\cap Y=V(\fp)$，其中 $\fp$ 是 $A$ 的素理想。可以验证 $V(\fp)=\ol{\{\fp\}}$。设 $y$ 是 $U$ 中对应于素理想 $\fp$ 的点。则 $\{y\}$ 在 $U$ 中的闭包为 $U\cap Y$。但 $U\cap Y$ 是不可约空间 $Y$ 的非空开子集，故它在 $Y$ 中稠密。于是 $Y=\ol{\{y\}}$。设 $y'$ 是另一个点，使 $Y=\ol{\{y'\}}$。则 $y'\in U\cap Y$。设 $\fp'$ 是 $A$ 中对应于 $y'$ 的素理想。则
\[
V(\fp')=\ol{\{\fp'\}}=U\cap Y=V(\fp).
\]
于是 $\fp=\fp'$。所以 $y=y'$。
\end{proof}

拓扑空间 $X$ 称为\emph{诺特}\idx{nuotekongjian}{诺特空间}的，若它满足以下等价条件之一：
\begin{enumerate}[label=(\roman*),leftmargin=2.8em]
  \item $X$ 的闭子集族满足降链条件，即对任意降链
  \[
  F_1\supset F_2\supset\cdots
  \]
  中的闭子集，当 $n$ 充分大时有 $F_n=F_{n+1}=\cdots$。
  \item $X$ 的开子集族满足升链条件。
  \item $X$ 的每个非空闭子集族都有极小元。
  \item $X$ 的每个非空开子集族都有极大元。
\end{enumerate}

可以证明，诺特拓扑空间的每个子空间都是诺特的。每个诺特拓扑空间 $X$ 都拟紧。事实上，给定开覆盖 $\{U_i\}_{i\in I}$，考虑 $X$ 的开子集族 $\bigcup_{i\in J}U_i$，其中 $J$ 取遍 $I$ 的有限子集。由 (iv)，此族有极大元。该极大元给出 $X$ 的有限开覆盖。

若 $A$ 是诺特环，则 $\Spec A$ 是诺特拓扑空间。事实上，给定 $\Spec A$ 中闭子集的降链
\[
V(\fa_1)\supset V(\fa_2)\supset\cdots,
\]
其中 $\fa_i$ 是 $A$ 的某些理想，我们就得到 $A$ 中理想的升链
\[
\sqrt{\fa_1}\subset \sqrt{\fa_2}\subset\cdots.
\]
因 $A$ 是诺特的，当 $n$ 充分大时有 $\sqrt{\fa_n}=\sqrt{\fa_{n+1}}=\cdots$。于是 $V(\fa_n)=V(\fa_{n+1})=\cdots$。

\begin{proposition}\label{it:1.3.4}
设 $X$ 是诺特拓扑空间。
\begin{enumerate}[label=(\roman*),leftmargin=2.8em]
  \item 对 $X$ 的每个闭子集 $Y$，有 $Y=Y_1\cup\cdots\cup Y_n$，其中 $Y_i$（$i=1,\dots,n$）是闭的不可约子集，且当 $i\ne j$ 时 $Y_i\not\subset Y_j$。这样的分解是唯一的。称 $Y_i$ 为 $Y$ 的\emph{不可约分支}\idx{bukeyuefenzhi}{不可约分支}。
  \item 不可约闭子集 $Y$ 是 $X$ 的不可约分支，当且仅当 $Y$ 在 $X$ 的不可约闭子集族中极大。
\end{enumerate}
\end{proposition}

\begin{proof}
(i) 设 $\cS$ 是 $X$ 的不能写成有限多个不可约闭子集之并的闭子集族。若 $\cS$ 非空，则因 $X$ 是诺特的，它有极小元 $Y$。因 $Y\in \cS$，$Y$ 不可约不成立。故 $Y=F_1\cup F_2$，其中 $F_1$ 与 $F_2$ 是 $Y$ 的真闭子集。由 $Y$ 的极小性，$F_1,F_2\notin \cS$。于是 $F_1$ 与 $F_2$ 都能写成有限多个不可约闭子集之并。但这样 $Y=F_1\cup F_2$ 也能写成有限多个不可约闭子集之并。这与 $Y\in \cS$ 矛盾。所以 $\cS$ 为空，即 $X$ 的每个闭子集都能写成有限多个不可约闭子集之并。

设
\[
Y=Y_1\cup\cdots\cup Y_m=Y_1'\cup\cdots\cup Y_n'
\]
是 $Y$ 的两个不可约闭子集分解，其中对所有 $i\ne j$ 有 $Y_i\not\subset Y_j$ 与 $Y_i'\not\subset Y_j'$。则 $Y_1=(Y_1\cap Y_1')\cup\cdots\cup (Y_1\cap Y_n')$。因 $Y_1$ 不可约，存在 $i$ 使 $Y_1=Y_1\cap Y_i'$，从而 $Y_1\subset Y_i'$。不妨设 $Y_1\subset Y_1'$。类似地，存在 $j$ 使 $Y_1'\subset Y_j$。因对任意 $j\ne 1$ 有 $Y_1\not\subset Y_j$，必有 $j=1$，从而 $Y_1=Y_1'$。类似地，重排 $Y_1',Y_2',\dots$ 的顺序，可得 $m=n$ 且 $Y_i=Y_i'$ 对每个 $i$ 成立。

(ii) 设 $X_1,\dots,X_n$ 是 $X$ 的不可约分支。对任意不可约闭子集 $Y$，有 $Y=(Y\cap X_1)\cup\cdots\cup (Y\cap X_n)$。故 $Y=Y\cap X_i$，从而 $Y\subset X_i$ 对某个 $i$ 成立。若 $Y$ 在不可约闭子集族中极大，则必有 $Y=X_i$，从而 $Y$ 是不可约分支。下面证明不可约分支 $X_1$ 在不可约闭子集族中极大。设 $X_1\subset Y$，其中 $Y$ 是不可约闭子集。如上所述，存在不可约分支 $X_i$ 使 $Y\subset X_i$。但这样 $X_1\subset X_i$。这蕴含 $i=1$ 且 $X_1=Y$。
\end{proof}

设 $A$ 是环。若 $A$ 的素理想不包含除自身外的任何素理想，则称它为\emph{极小素理想}\idx{jixiaosulixiang}{极小素理想}。由 Zorn 引理可以证明，$A$ 的每个素理想都包含一个极小素理想。

\begin{corollary}\label{it:1.3.5}
设 $A$ 是环，使 $\Spec A$ 是诺特拓扑空间。则 $\Spec A$ 的不可约分支族与 $A$ 的极小素理想族之间存在一一对应。
\end{corollary}

\begin{proof}
由 \ref{it:1.3.2} 与 \ref{it:1.2.2}(iv)，映射 $\fp\mapsto V(\fp)$ 在 $A$ 的素理想族与 $\Spec A$ 的不可约闭子集族之间建立一一对应。由 \ref{it:1.3.4}(ii)，$\Spec A$ 的每个不可约分支都是极大的不可约闭子集。若把它写成 $V(\fp)$ 的形式，其中 $\fp$ 是素理想，则 $\fp$ 是极小素理想。
\end{proof}

概形 $(X,\cO_X)$ 称为\emph{连通}\idx{liantonggaixing}{连通概形}（resp.\ \emph{不可约}）的，若 $X$ 连通（resp.\ 不可约）。若对 $X$ 的每个开子集 $U$，$\cO_X(U)$ 是既约环（resp.\ 整环），则称它为\emph{既约}（resp.\ \emph{整}）的。

\begin{proposition}\label{it:1.3.6}
$(X,\cO_X)$ 是整概形，当且仅当它不可约且既约。
\end{proposition}

\begin{proof}
设 $(X,\cO_X)$ 是整概形。由定义它是既约的。若它不可约不成立，则可找到 $X$ 的两个不交非空开子集 $U_1$ 与 $U_2$。有 $\cO_X(U_1\cup U_2)=\cO_X(U_1)\times \cO_X(U_2)$，它不是整环。（有 $(1,0)(0,1)=(0,0)$，但 $(1,0),(0,1)\ne(0,0)$。）矛盾。所以 $X$ 不可约。

反之，设 $(X,\cO_X)$ 不可约且既约。设 $U$ 是 $X$ 的开子集，$f_1,f_2\in \cO_X(U)$ 满足 $f_1f_2=0$。对每个 $i\in\{1,2\}$，设 $U_{f_i}$ 为 $U$ 中满足「$f_i$ 在 $P$ 处的芽是局部环 $\cO_{X,P}$ 中的单位」的点 $P$ 构成的子集。用仿射开子概形 $U_\alpha=\Spec A_\alpha$（$\alpha\in I$）覆盖 $U$，并把限制 $f_i|_{U_\alpha}\in \cO_X(U_\alpha)\cong A_\alpha$ 记为 $f_{i\alpha}$。则对所有 $\alpha$ 有 $U_{f_i}\cap U_\alpha=D(f_{i\alpha})$。故每个 $U_{f_i}$ 是开集。此外 $U_{f_1}\cap U_{f_2}=U_{f_1f_2}$。因 $f_1f_2=0$，有 $U_{f_1f_2}=\emptyset$。故 $U_{f_1}\cap U_{f_2}=\emptyset$。若 $U$ 不可约，则 $U_{f_1}$ 或 $U_{f_2}$ 为空，不妨设 $U_{f_1}=\emptyset$。则对每个 $\alpha$ 有 $D(f_{1\alpha})=U_{f_1}\cap U_\alpha=\emptyset$。这蕴含 $f_{1\alpha}$ 幂零。因 $\cO_X(U_\alpha)$ 既约，对每个 $\alpha$ 有 $f_{1\alpha}=0$。于是 $f_1=0$。这说明 $\cO_X(U)$ 是整环。
\end{proof}

设 $X$ 是整概形。设 $\xi$ 是它的一般点。则 $\cO_{X,\xi}$ 是域，我们称它为 $X$ 的\emph{函数域}\idx{hanshuyu}{函数域}。事实上，对 $X$ 的每个仿射开子概形 $U=\Spec A$，$A$ 必是整环，$\xi$ 在 $U$ 中，且 $A$ 中对应于 $\xi$ 的素理想 $\fp$ 为 $(0)$。于是 $\cO_{X,\xi}$ 同构于 $A$ 的分式域 $A_\fp$。这说明 $X$ 的函数域同构于 $\cO_X(U)$ 的分式域，其中 $U$ 是 $X$ 的任意仿射开子概形。对每个 $P\in X$ 与 $P$ 的每个开邻域 $U$，有 $\xi\in U$。环同态族 $\cO_X(U)\to \cO_{X,\xi}$ 诱导同态
\[
\cO_{X,P}=\dlim_{P\in U}\cO_X(U)\to \cO_{X,\xi}.
\]
此同态是单射，且对每个 $P\in X$，它把 $\cO_{X,\xi}$ 与 $\cO_{X,P}$ 的分式域等同起来。对开子集的每个包含 $U\subset V$ 与 $U$ 中的每个 $P$，典范同态
\[
\cO_X(V)\to \cO_X(U)\to \cO_{X,P}\to \cO_{X,\xi}
\]
都是单射。

下面的命题留给读者证明：
\begin{proposition}\label{it:1.3.7}
概形 $(X,\cO_X)$ 既约，当且仅当对每个 $P\in X$，$\cO_{X,P}$ 既约。
\end{proposition}

概形 $(X,\cO_X)$ 称为\emph{局部整}的，若对每个 $P\in X$，$\cO_{X,P}$ 是整环。每个整概形都是局部整的。

\begin{proposition}\label{it:1.3.8}
设 $(X,\cO_X)$ 是概形，使 $X$ 是诺特拓扑空间。则 $(X,\cO_X)$ 局部整，当且仅当它既约且其不可约分支两两不交。任意局部整诺特概形的连通分支都与不可约分支重合。
\end{proposition}

\begin{proof}
设 $X$ 的不可约分支 $X_i$（$i=1,\dots,n$）两两不交。则每个 $X_i$ 既开又闭。若 $(X,\cO_X)$ 既约，则每个 $(X_i,\cO_{X_i})$ 既既约又不可约。故每个 $(X_i,\cO_{X_i})$ 是整概形。于是对每个 $P\in X_i$，$\cO_{X,P}$ 是整环。所以 $(X,\cO_X)$ 局部整。

反之，设 $(X,\cO_X)$ 局部整。则由 \ref{it:1.3.7} 它既约。设 $X_1$ 与 $X_2$ 是 $X$ 的两个不可约分支，$P\in X_1\cap X_2$。设 $U=\Spec A$ 是 $X$ 的包含 $P$ 的仿射开子概形。可以证明 $X_1\cap U$ 与 $X_2\cap U$ 都是 $U$ 的不可约分支。由 \ref{it:1.3.5}，有 $X_i\cap U=V(\fp_i)$（$i=1,2$），其中 $\fp_i$ 是 $A$ 的极小素理想。设 $\fp$ 是 $A$ 中对应于点 $P\in X_1\cap X_2\cap U$ 的素理想。有 $\fp_i\subset \fp$。故 $\fp_1A_\fp$ 与 $\fp_2A_\fp$ 是 $A_\fp$ 的两个极小素理想。因 $A_\fp=\cO_{X,P}$ 是整环，它只有一个极小素理想。故 $\fp_1=\fp_2$，于是 $V(\fp_1)=V(\fp_2)$，即 $X_1\cap U=X_2\cap U$。但对每个 $i\in\{1,2\}$，$X_i\cap U$ 在 $X_i$ 中稠密。于是 $X_1=X_2$。所以 $X$ 的不可约分支两两不交。
\end{proof}

设 $(X,\cO_X)$ 是概形。对任意截面 $f\in \cO_X(X)$，定义 $X_f$ 为 $X$ 中满足「$f$ 在 $P$ 处的芽是 $\cO_{X,P}$ 中的单位」的点 $P$ 构成的子集。

\begin{proposition}\label{it:1.3.9}
设 $(X,\cO_X)$ 是概形。
\begin{enumerate}[label=(\roman*),leftmargin=2.8em]
  \item 对每个 $f\in \cO_X(X)$，$X_f$ 是开集。$X_f$ 为空，当且仅当存在 $X$ 的开覆盖 $\{U_i\}_{i\in I}$，使每个 $f|_{U_i}$ 幂零。对任意 $f,g\in \cO_X(X)$，有 $X_f\cap X_g=X_{fg}$。
  \item 设 $(\varphi,\varphi^\sharp)\colon (X,\cO_X)\to (Y,\cO_Y)$ 是概形的态射，$f\in \cO_Y(Y)$，并设 $\varphi^\sharp(f)$ 是 $f$ 在同态 $\varphi^\sharp(Y)\colon \cO_Y(Y)\to \cO_X(X)$ 下的像。则 $\varphi^{-1}(Y_f)=X_{\varphi^\sharp(f)}$。
  \item 设 $X$ 可被有限多个仿射开子概形 $U_i$（$i\in I$）覆盖，且对所有 $i,j\in I$，$U_i\cap U_j$ 可被有限多个仿射开子概形覆盖。设 $A=\cO_X(X)$。则对每个 $f\in A$，有 $\cO_X(X_f)\cong A_f$。
\end{enumerate}
\end{proposition}

\begin{tnote}
仿射概形情形下的主开集 $D(f)$，在一般概形上自然推广为 $X_f$；下面的 \ref{it:1.3.10} 正是用它给出仿射性的判别。[依据：原书 p.35–36]
\end{tnote}

\begin{proof}
(i) 参看 \ref{it:1.3.6} 的证明。

(ii) 对每个 $P\in X$，$f$ 在 $\varphi(P)$ 处的芽 $f_{\varphi(P)}$ 在同态 $\varphi^\sharp_P\colon \cO_{Y,\varphi(P)}\to \cO_{X,P}$ 下的像，就是 $\varphi^\sharp(f)$ 在 $P$ 处的芽 $(\varphi^\sharp(f))_P$。因该同态是局部的，$f_{\varphi(P)}$ 是单位当且仅当 $(\varphi^\sharp(f))_P$ 是单位，即 $\varphi(P)\in Y_f$ 当且仅当 $P\in X_{\varphi^\sharp(f)}$。

(iii) 限制 $A=\cO_X(X)\to \cO_X(X_f)$ 把 $f$ 送到单位。故它诱导同态 $A_f\to \cO_X(X_f)$。下面证明它是同构。

对每个 $i$，设 $U_i=\Spec A_i$。把 $f$ 在限制 $\cO_X(X)\to \cO_X(U_i)\cong A_i$ 下的像记为 $f_i$。则 $X_f\cap U_i$ 与 $U_i$ 的开子集 $D(f_i)$ 重合。

设 $s\in \cO_X(X)$ 是使 $s|_{X_f}=0$ 的截面。我们证明存在自然数 $n$，使 $f^n s=0$ 在 $\cO_X(X)$ 中成立。这将证明同态 $A_f\to \cO_X(X_f)$ 是单射。

注意对每个 $i$，$s|_{U_i}$ 属于限制 $\cO_X(U_i)\to \cO_X(U_i\cap X_f)$ 的核。此限制可与典范同态 $A_i\to (A_i)_{f_i}$ 等同。因 $A_i\to (A_i)_{f_i}$ 的核中每个元素都被 $f_i$ 的某个幂零化，存在自然数 $n_i$，使 $f_i^{n_i}s|_{U_i}=0$ 在 $\cO_X(U_i)$ 中成立。因 $\{U_i\}_{i\in I}$ 是有限覆盖，可取充分大的 $n$，使对每个 $i$ 有 $f^n s|_{U_i}=0$。于是 $f^n s=0$。

其次证明：对任意截面 $t\in \cO_X(X_f)$，存在自然数 $n$，使 $f^n t$ 可延拓为 $\cO_X(X)=A$ 中的截面。这将证明同态 $A_f\to \cO_X(X_f)$ 是满射。

因限制 $\cO_X(U_i)\to \cO_X(U_i\cap X_f)$ 可与典范同态 $A_i\to (A_i)_{f_i}$ 等同，存在自然数 $n$，使对每个 $i$，$f_i^n t|_{U_i\cap X_f}$ 可延拓为 $\cO(U_i)$ 中的截面 $t_i$。注意 $t_i|_{U_i\cap U_j}-t_j|_{U_i\cap U_j}$ 在 $U_i\cap U_j$ 上的限制，在 $U_i\cap U_j\cap X_f$ 上为零，其中 $i,j$ 任意。由上面所证，存在自然数 $m$，使对所有 $i,j\in I$ 有 $f^m(t_i|_{U_i\cap U_j}-t_j|_{U_i\cap U_j})=0$。于是可把 $f_i^m t_i$ 粘起来，得到 $\cO_X(X)$ 中的截面。该截面是 $f^{m+n}t$ 的延拓。
\end{proof}

\begin{proposition}\label{it:1.3.10}
概形 $(X,\cO_X)$ 是仿射概形，当且仅当存在有限多个截面 $f_1,\dots,f_n\in \cO(X)$，它们生成 $\cO(X)$ 的单位理想，且使每个开子概形 $(X_{f_i},\cO_X|_{X_{f_i}})$ 都是仿射的。
\end{proposition}

\begin{proof}
先证「当」的部分。设 $A=\cO(X)$。由 \ref{it:1.2.8}，恒等同态 $A\to \cO(X)$ 定义态射 $\varphi\colon X\to \Spec A$。下面证明 $\varphi$ 是同构。因 $f_i$（$i=1,\dots,n$）生成单位理想，$D(f_i)$（$i=1,\dots,n$）构成 $\Spec A$ 的开覆盖。由 \ref{it:1.3.9}(ii) 有 $\varphi^{-1}(D(f_i))=X_{f_i}$。我们验证 \ref{it:1.3.9}(iii) 的条件成立，从而 $\cO_X(X_{f_i})\cong A_{f_i}$。每个 $X_{f_i}$ 由假设是仿射的。故由 \ref{it:1.2.5}，$\varphi$ 诱导同构 $\varphi^{-1}(D(f_i))\cong D(f_i)$，对每个 $i$ 成立，从而 $\varphi$ 是同构。

下面验证 \ref{it:1.3.9}(iii) 的条件成立。$X$ 可被 $X_{f_i}$（$i=1,\dots,n$）覆盖，且每个 $X_{f_i}$ 是仿射的。设 $X_{f_i}=\Spec B_i$。对每个 $j$，设 $f_j'$ 是 $f_j$ 在限制 $\cO_X(X)\to \cO_X(X_{f_i})=B_i$ 下的像。由 \ref{it:1.3.9}(ii)，$X_{f_i}\cap X_{f_j}$ 等于 $\Spec B_i$ 的子集 $D(f_j')$，它是仿射的。
\end{proof}

概形 $(X,\cO_X)$ 称为\emph{局部诺特}的，若它可被仿射开子概形 $U_i=\Spec A_i$（$i\in I$）覆盖，且每个 $A_i$ 是诺特环。若可找到 $X$ 的有限仿射开覆盖 $U_i=\Spec A_i$（$i\in I$），使每个 $A_i$ 是诺特环，则称它为\emph{诺特}的。$(X,\cO_X)$ 是诺特概形，当且仅当它局部诺特且拟紧。若 $(X,\cO_X)$ 是诺特概形，则 $X$ 是诺特拓扑空间。但反过来不成立。

\begin{proposition}\label{it:1.3.11}
设 $(X,\cO_X)$ 是局部诺特概形。则对 $X$ 的任意仿射开子概形 $U=\Spec A$，$A$ 是诺特环。特别地，仿射概形 $(\Spec A,\cO_{\Spec A})$ 局部诺特，当且仅当 $A$ 是诺特环。
\end{proposition}

\begin{proof}
用仿射开子概形 $U_i=\Spec A_i$（$i\in I$）覆盖 $X$，使每个 $A_i$ 是诺特的。则 $U=\Spec A$ 可被 $U_i\cap U$（$i\in I$）覆盖。因每个 $U_i\cap U$ 在 $U_i=\Spec A_i$ 中是开的，我们可以用 $D(f_{ij})\cong \Spec (A_i)_{f_{ij}}$（$j\in J_i$）覆盖 $U_i\cap U$，其中 $f_{ij}\in A_i$。注意每个 $(A_i)_{f_{ij}}$ 是诺特的。改变记号，我们可以用仿射开子概形 $U_i=\Spec A_i$（$i\in I$）覆盖 $U=\Spec A$，使每个 $A_i$ 是诺特的。用 $D(f_{ij})$（$j\in J_i$）覆盖每个 $U_i$，其中 $f_{ij}\in A$。对每个 $j\in J_i$，设 $f_{ij}'$ 是 $f_{ij}$ 在限制 $\cO(U)\cong A\to \cO(U_i)\cong A_i$ 下的像。则由 \ref{it:1.3.9}(ii)，$\Spec A$ 的子集 $D(f_{ij})$ 与 $\Spec A_i$ 的子集 $D(f_{ij}')$ 重合。于是 $\cO(D(f_{ij}))=\cO(D(f_{ij}'))$，从而 $A_{f_{ij}}\cong (A_i)_{f_{ij}'}$。因 $A_i$ 是诺特的，$(A_i)_{f_{ij}'}$ 是诺特的，从而 $A_{f_{ij}}$ 是诺特的。改变记号，我们可以用开子集 $D(f_i)$（$i\in I$）覆盖 $U=\Spec A$，其中 $f_i\in A$，且每个 $A_{f_i}$ 是诺特的。因 $\Spec A$ 拟紧，可设此覆盖是有限的。对 $A$ 的任意理想 $\fa$，设 $\{\frac{a_{ij}}{f_i^{k_{ij}}}\}$（$j\in J_i$）是 $\fa A_{f_i}$ 的有限生成元族。我们断言 $\fa$ 由有限族 $\{a_{ij}\}$ 生成。给定 $a\in \fa$，对每个 $i$，可把
\[
\frac{a}{1}=\sum_j \frac{b_{ij}}{f_i^{k_{ij}}}\frac{a_{ij}}{f_i^{k_{ij}}}
\]
写成 $A_{f_i}$ 中的等式。故对充分大的 $k_i$，在 $A$ 中有
\[
f_i^{k_i}a=\sum_j c_{ij}a_{ij}.
\]
但 $\bigcup_i D(f_i^{k_i})=\Spec A$。故 $V(\sum(f_i^{k_i}))=V(A)$，从而 $1\in \sqrt{\sum(f_i^{k_i})}$。设 $1=\sum_i b_if_i^{k_i}$。则
\[
a=\sum_i b_if_i^{k_i}a=\sum_{i,j} b_ic_{ij}a_{ij}.
\]
于是 $\fa$ 由有限族 $\{a_{ij}\}$ 生成。所以 $A$ 是诺特的。
\end{proof}

概形的态射 $f\colon X\to Y$ 称为\emph{拟紧}\idx{nijintaishe}{拟紧态射}的，若存在 $Y$ 的仿射开子概形覆盖 $V_i$（$i\in I$），使每个 $f^{-1}(V_i)$ 拟紧。若 $Y$ 可被仿射开子概形 $V_i$（$i\in I$）覆盖，使每个 $f^{-1}(V_i)$ 是仿射的，则称 $f$ 为\emph{仿射}态射。若 $Y$ 可被仿射开子概形 $V_i=\Spec B_i$（$i\in I$）覆盖，使每个 $f^{-1}(V_i)$ 可被仿射开子概形 $U_{ij}=\Spec A_{ij}$（$j\in J_i$）覆盖，其中 $A_{ij}$ 是有限生成的 $B_i$-代数，则称 $f$ \emph{局部有限型}。若 $f$ 既拟紧又局部有限型，则称它\emph{有限型}。若 $Y$ 可被仿射开子概形 $V_i=\Spec B_i$（$i\in I$）覆盖，使对每个 $i$ 有 $f^{-1}(V_i)=\Spec A_i$，其中 $A_i$ 是某个作为 $B_i$-模有限生成的 $B_i$-代数，则称 $f$ \emph{有限}。

\begin{proposition}\label{it:1.3.12}
设 $f\colon X\to Y$ 是概形的态射。
\begin{enumerate}[label=(\roman*),leftmargin=2.8em]
  \item $f$ 拟紧，当且仅当对 $Y$ 的每个拟紧开子集 $V$，$f^{-1}(V)$ 拟紧。
  \item $f$ 仿射，当且仅当对 $Y$ 的每个仿射开子概形 $V$，$f^{-1}(V)$ 是仿射的。
  \item $f$ 局部有限型，当且仅当对 $Y$ 的每个仿射开子概形 $V=\Spec B$ 与 $X$ 的每个满足 $f(U)\subset V$ 的仿射开子概形 $U=\Spec A$，$B$-代数 $A$ 是有限生成的。
  \item[(v)] $f$ 有限型，当且仅当对 $Y$ 的每个仿射开子概形 $V=\Spec B$，$f^{-1}(V)$ 可被有限多个仿射开子概形 $U_j=\Spec A_j$（$j\in J$）覆盖，使每个 $A_j$ 是有限生成的 $B$-代数。
  \item[(vi)] $f$ 有限，当且仅当对 $Y$ 的每个仿射开子概形 $V=\Spec B$，$f^{-1}(V)=\Spec A$，其中 $A$ 是某个作为 $B$-模有限生成的 $B$-代数。
\end{enumerate}
\end{proposition}

\begin{tnote}
原书 \ref{it:1.3.12} 的分条编号从 (iii) 直接跳到 (v)，没有 (iv)；译文保留原编号，以免与原文对照时错位。[依据：原书 p.38]
\end{tnote}

\begin{proof}
我们只证 (ii)，其余留给读者。设 $f\colon X\to Y$ 是仿射态射。用仿射开子概形 $V_i=\Spec A_i$（$i\in I$）覆盖 $Y$，使对某个 $A_i$-代数 $B_i$ 有 $f^{-1}(V_i)=\Spec B_i$。对任意仿射开子概形 $V=\Spec A$，用形如 $D(f_{ij})$（$j\in J_i$）的开子集覆盖每个 $V_i\cap V$，其中 $f_{ij}\in A_i$。注意 $D(f_{ij})$（$i\in I$，$j\in J_i$）构成 $V$ 的开覆盖，每个 $D(f_{ij})$ 是仿射的，且 $f^{-1}(D(f_{ij}))\cong \Spec(B_i)_{f_{ij}}$ 也是仿射的。改变记号，可设 $V=\Spec A$ 可被仿射开子概形 $V_i=\Spec A_i$（$i\in I$）覆盖，使对某个 $A_i$-代数 $B_i$ 有 $f^{-1}(V_i)=\Spec B_i$。用形如 $D(f_{ij})$（$j\in J_i$）的开子集覆盖每个 $V_i$，其中 $f_{ij}\in A$。对每个 $i$ 与每个 $j\in J_i$，设 $f_{ij}'\in A_i$ 是 $f_{ij}$ 在限制 $A\cong \cO_Y(V)\to A_i\cong \cO_Y(V_i)$ 下的像。$V=\Spec A$ 的子集 $D(f_{ij})$ 与 $V_i=\Spec A_i$ 的子集 $D(f_{ij}')$ 重合。因 $f^{-1}(D(f_{ij}'))\cong \Spec(B_i)_{f_{ij}'}$，每个 $f^{-1}(D(f_{ij}))$ 是仿射的。改变记号，可设 $V=\Spec A$ 可被 $D(f_i)$（$i\in I$）覆盖，其中 $f_i\in A$，且每个 $f^{-1}(D(f_i))$ 是仿射的。因 $\Spec A$ 拟紧，可设此覆盖是有限的。可以验证 $f_i$（$i\in I$）生成 $A$ 的单位理想。把 $f^{-1}(V)$ 记为 $U$，并设 $f_i'$（$i\in I$）是 $f_i$ 在同态 $A\cong \cO_Y(V)\to \cO_X(U)$ 下的像。则 $f_i'$（$i\in I$）生成 $\cO_X(U)$ 的单位理想。此外由 \ref{it:1.3.9}(ii) 有 $U_{f_i'}=f^{-1}(D(f_i))$。故每个 $U_{f_i'}$ 是仿射的。由 \ref{it:1.3.10}，$U=f^{-1}(V)$ 是仿射的。
\end{proof}

\begin{proposition}\label{it:1.3.13}
设 $X\to S$ 与 $Y\to S$ 是两个态射，并设 $Y\to S$ 局部有限型。
\begin{enumerate}[label=(\roman*),leftmargin=2.8em]
  \item 设 $f,g\colon X\to Y$ 是两个态射，使下图交换：
  \begin{center}
  \begin{tikzcd}
  X \arrow[r] \arrow[dr] & Y \arrow[d] \\
  & S
  \end{tikzcd}
  \end{center}
  设 $x\in X$，使 $f(x)=g(x)=y$，且 $f^\sharp_x,g^\sharp_x\colon \cO_{Y,y}\to \cO_{X,x}$ 相同。则存在 $x$ 在 $X$ 中的开邻域 $U$，使 $f|_U=g|_U$。
  \item 设 $S$ 局部诺特。设 $x\in X$ 与 $y\in Y$ 是两点，使它们在 $S$ 中的像为同一点 $s\in S$，并设 $\phi\colon \cO_{Y,y}\to \cO_{X,x}$ 是使下图交换的同态：
  \begin{center}
  \begin{tikzcd}
  \cO_{X,x} & \cO_{Y,y} \arrow[l, "\phi"'] \\
  & \cO_{S,s} \arrow[ul] \arrow[u]
  \end{tikzcd}
  \end{center}
  则存在 $x$ 在 $X$ 中的开邻域 $U$ 与态射 $f\colon U\to Y$，使 $f(x)=y$，$f^\sharp_x=\phi$，且下图交换：
  \begin{center}
  \begin{tikzcd}
  U \arrow[r, "f"] \arrow[dr] & Y \arrow[d] \\
  & S
  \end{tikzcd}
  \end{center}
  \item 设 $S$ 局部诺特，$X\to S$ 也局部有限型，$f\colon X\to Y$ 是使 (i) 中图表交换的态射。设 $f(x)=y$，且 $f^\sharp_x\colon \cO_{Y,y}\to \cO_{X,x}$ 是同构。则存在 $x$ 在 $X$ 中的开邻域 $U$ 与 $y$ 在 $Y$ 中的开邻域 $V$，使 $f$ 诱导从 $U$ 到 $V$ 的同构。
\end{enumerate}
\end{proposition}

\begin{tnote}
\ref{it:1.3.13} 的 (i)(ii) 表明，局部有限型态射在一点附近的性状由该点处的局部同态决定；这正是 \ref{it:1.3.14} 中双有理态射局部为同构的依据。[依据：原书 p.39–42]
\end{tnote}

\begin{proof}
(i) 设 $s$ 是 $y$ 在 $S$ 中的像，$W=\Spec C$ 是 $s$ 在 $S$ 中的仿射开邻域，$V=\Spec B$ 是 $y$ 在 $Y$ 中的仿射开邻域，它在态射 $Y\to S$ 下映入 $W$，且 $B$ 是有限生成的 $C$-代数，并设 $U=\Spec A$ 是 $x$ 在 $X$ 中的仿射开邻域，使 $f(U),g(U)\subset V$。设 $\fp,\fq,\fr$ 分别是 $A,B,C$ 中对应于 $x,y,s$ 的素理想。则 $f$ 与 $g$ 分别定义 $C$-代数同态 $\phi,\psi\colon B\to A$，使 $\phi^{-1}(\fp)=\psi^{-1}(\fp)=\fq$，且 $\phi$ 与 $\psi$ 诱导相同的同态 $B_\fq\to A_\fp$。设 $b_1,\dots,b_n$ 是 $C$-代数 $B$ 的有限生成元族。则在 $A_\fp$ 中有 $\frac{\phi(b_i)}{1}=\frac{\psi(b_i)}{1}$。故对每个 $i$，存在 $s_i\in A-\fp$，使 $s_i\phi(b_i)=s_i\psi(b_i)$。设 $s=s_1\cdots s_n$。则在 $A_s$ 中有 $\frac{\phi(b_i)}{1}=\frac{\psi(b_i)}{1}$，从而 $\phi$ 与 $\psi$ 诱导相同的同态 $B\to A_s$。这说明 $f$ 与 $g$ 在 $x$ 的开邻域 $D(s)$ 上的限制相同。

(ii) 设 $W=\Spec C$ 是 $s$ 在 $S$ 中的仿射开邻域，$V=\Spec B$ 是 $y$ 在 $Y$ 中的仿射开邻域，它在态射 $Y\to S$ 下映入 $W$，且 $B$ 是有限生成的 $C$-代数，并设 $U=\Spec A$ 是 $x$ 在 $X$ 中的仿射开邻域，它在态射 $X\to S$ 下映入 $W$。设 $\fp$ 与 $\fq$ 分别是 $A$ 与 $B$ 中对应于 $x$ 与 $y$ 的素理想。则 $\phi\colon \cO_{Y,y}\to \cO_{X,x}$ 对应于 $C$-代数同态 $B_\fq\to A_\fp$，我们也把它记为 $\phi$。设 $b_1,\dots,b_n$ 是 $C$-代数 $B$ 的有限生成元族。则我们有 $C$-代数的满同态
\[
C[x_1,\dots,x_n]\to B,\quad x_i\mapsto b_i.
\]
设 $I$ 是它的核。因 $C[x_1,\dots,x_n]$ 是诺特的，$I$ 是有限生成的。设
$f_1(x_1,\dots,x_n)$, $\dots$, $f_m(x_1,\dots,x_n)$ 是理想 $I$ 在 $C[x_1,\dots,x_n]$ 中的有限生成元族。取充分大的 $d$，可把每个 $f_i(x_1,\dots,x_n)$ 写成
\[
f_i(x_1,\dots,x_n)=\sum_{0\le i_1,\dots,i_n\le d} c_{i i_1\dots i_n}x_1^{i_1}\cdots x_n^{i_n}.
\]
把复合
\[
C[x_1,\dots,x_n]\to B\to B_\fq\xrightarrow{\phi} A_\fp
\]
记为 $\theta$。对某个 $a_i\in A$ 与 $s_i\in A-\fp$，有 $\phi(\frac{b_i}{1})=\frac{a_i}{s_i}$。在 $A_\fp$ 中有
\[
\theta(f_i(x_1,\dots,x_n))=0,
\]
即
\[
\phi\left(\frac{f_i(b_1,\dots,b_n)}{1}\right)=0.
\]
因
\[
\begin{aligned}
\phi\left(\frac{f_i(b_1,\dots,b_n)}{1}\right)
&=\sum_{0\le i_1,\dots,i_n\le d} c_{i i_1\dots i_n}\left(\frac{a_1}{s_1}\right)^{i_1}\cdots\left(\frac{a_n}{s_n}\right)^{i_n}\\
&=\frac{\sum_{0\le i_1,\dots,i_n\le d} c_{i i_1\dots i_n}s_1^{d-i_1}a_1^{i_1}\cdots s_n^{d-i_n}a_n^{i_n}}{s_1^d\cdots s_n^d},
\end{aligned}
\]
则存在元素 $t_i\in A-\fp$，使
\[
t_i\sum_{0\le i_1,\dots,i_n\le d} c_{i i_1\dots i_n}s_1^{d-i_1}a_1^{i_1}\cdots s_n^{d-i_n}a_n^{i_n}=0.
\]
设 $s=s_1\cdots s_n t_1\cdots t_m$，并考虑 $C$-代数同态
\[
C[x_1,\dots,x_n]\to A_s,\quad x_i\mapsto \frac{a_i}{s_i}=\frac{a_i s_1\cdots\hat s_i\cdots s_n t_1\cdots t_m}{s}.
\]
它把每个 $f_i(x_1,\dots,x_n)$ 送到 $0$，故 $I$ 包含在它的核中。于是它诱导同态 $B\to A_s$，且图表
\begin{center}
\begin{tikzcd}
B \arrow[r] \arrow[d] & A_s \arrow[d] \\
B_\fq \arrow[r, "\phi"] & A_\fp
\end{tikzcd}
\end{center}
交换，其中竖箭头是典范同态。同态 $B\to A_s$ 诱导概形的态射 $f\colon D(s)\to Y$，具有所需性质。

(iii) 由 (ii)，存在 $y$ 在 $Y$ 中的邻域 $V$ 与态射 $g\colon V\to X$，使 $g(y)=x$，$g^\sharp_y={f^\sharp_x}^{-1}$，且图表
\begin{center}
\begin{tikzcd}
V \arrow[r, "g"] \arrow[dr] & X \arrow[d] \\
& S
\end{tikzcd}
\end{center}
交换。由 (i)，存在 $x$ 在 $f^{-1}(V)$ 中的邻域 $U$，使复合
\[
U\xrightarrow{f}V\xrightarrow{g}X
\]
与典范态射 $U\hookrightarrow X$ 重合；且存在 $y$ 在 $V$ 中的邻域 $V'$，使复合
\[
V'\xrightarrow{g}X\xrightarrow{f}Y
\]
与典范态射 $V'\hookrightarrow Y$ 重合。我们断言 $f$ 与 $g$ 诱导同构
\[
U\cap f^{-1}(V')\xrightarrow{f}g^{-1}(U)\cap V',
\qquad
g^{-1}(U)\cap V'\xrightarrow{g}U\cap f^{-1}(V'),
\]
它们互为逆。事实上，对任意 $P\in U\cap f^{-1}(V')$，有 $f(P)\in V'$。此外，因 $P\in U$，有 $gf(P)=P$，从而 $f(P)\in g^{-1}(U)$。故 $f(P)\in g^{-1}(U)\cap V'$。于是
\[
f(U\cap f^{-1}(V'))\subset g^{-1}(U)\cap V',
\]
且 $f$ 诱导态射 $U\cap f^{-1}(V')\to g^{-1}(U)\cap V'$。类似地，
\[
g(g^{-1}(U)\cap V')\subset U\cap f^{-1}(V'),
\]
从而 $g$ 诱导态射 $g^{-1}(U)\cap V'\to U\cap f^{-1}(V')$。此外有
\[
(gf)|_{U\cap f^{-1}(V')}=\id_{U\cap f^{-1}(V')},\qquad
(fg)|_{g^{-1}(U)\cap V'}=\id_{g^{-1}(U)\cap V'}.
\]
\end{proof}

设 $X$ 与 $Y$ 是两个整概形。若 $\ol{f(X)}=Y$，则称态射 $f\colon X\to Y$ 为\emph{支配}的。设 $\xi$ 是 $X$ 的一般点，$\eta$ 是 $Y$ 的一般点。有
\[
\ol{\{f(\xi)\}}=f(\ol{\{\xi\}})=\ol{f(X)}=Y=\ol{\{\eta\}}.
\]
故必有 $f(\xi)=\eta$。若函数域上的同态 $f^\sharp_\xi\colon \cO_{Y,\eta}\to \cO_{X,\xi}$ 是同构，则称 $f$ 为\emph{双有理}的。

\begin{corollary}\label{it:1.3.14}
设 $S$ 是局部诺特概形，$X$ 与 $Y$ 是两个整概形。假设有交换图表
\begin{center}
\begin{tikzcd}
X \arrow[rr, "f"] \arrow[dr] & & Y \arrow[dl] \\
& S
\end{tikzcd}
\end{center}
使 $X\to S$ 与 $Y\to S$ 局部有限型，且 $f$ 是双有理态射。则存在非空开子集 $U\subset X$ 与 $V\subset Y$，使 $f$ 诱导从 $U$ 到 $V$ 的同构。
\end{corollary}

\begin{proof}
设 $\xi$ 是 $X$ 的一般点，$\eta$ 是 $Y$ 的一般点。有 $f(\xi)=\eta$，且 $f^\sharp_\xi\colon \cO_{Y,\eta}\to \cO_{X,\xi}$ 是同构。于是我们的断言由 \ref{it:1.3.13}(iii) 得到。
\end{proof}

设 $(X,\cO_X)$ 是概形。回忆 $(X,\cO_X)$ 的开子概形具有形式 $(U,\cO_X|_U)$，其中 $U$ 是 $X$ 的开子集。若态射 $(f,f^\sharp)\colon (Z,\cO_Z)\to (X,\cO_X)$ 诱导 $(Z,\cO_Z)$ 与 $(X,\cO_X)$ 的某个开子概形的同构，即 $f$ 诱导 $Z$ 与 $X$ 的某个开子集的同胚，且对 $X$ 的每个包含在 $f(Z)$ 中的开子集 $U$，$f^\sharp(U)\colon \cO_X(U)\to \cO_Z(f^{-1}(U))$ 是同构，则称它为\emph{开浸入}\idx{kaijinru}{开浸入}。

若 $f\colon Z\to X$ 诱导 $Z$ 与 $X$ 的某个闭子集的同胚，且 $f^\sharp\colon \cO_X\to f_*\cO_Z$ 是满射，则称态射 $(f,f^\sharp)\colon (Z,\cO_Z)\to (X,\cO_X)$ 为\emph{闭浸入}\idx{bijinru}{闭浸入}。若存在同构（必唯一）$Z_1\to Z_2$，使图表
\begin{center}
\begin{tikzcd}
Z_1 \arrow[r] & X \\
Z_2 \arrow[u] \arrow[ur] &
\end{tikzcd}
\end{center}
交换，则称两个闭浸入 $Z_1\to X$ 与 $Z_2\to X$ 同构。$X$ 的闭浸入的同构类称为 $X$ 的\emph{闭子概形}\idx{bizigaixing}{闭子概形}。

若态射 $Z\to X$ 可写成复合 $Z\to U\to X$，其中 $U\to X$ 是开浸入且 $Z\to U$ 是闭浸入，则称它为\emph{浸入}。若 $X$ 的子集 $Z$ 是开子集与闭子集的交，则称它为\emph{局部闭集}\idx{jububiji}{局部闭集}。这等价于说 $Z$ 在 $\ol{Z}$ 中是开的。若 $Z\to X$ 是概形的浸入，则 $Z$ 在 $X$ 中的像是局部闭子集。

下面命题留给读者证明：

\begin{proposition}\label{it:1.3.15}
设 $(f,f^\sharp)\colon (Z,\cO_Z)\to (X,\cO_X)$ 是态射。
\begin{enumerate}[label=(\roman*),leftmargin=2.8em]
  \item $(f,f^\sharp)$ 是开浸入，当且仅当 $f$ 诱导 $Z$ 与 $X$ 的某个开子集的同胚，且对每个 $P\in Z$，$f^\sharp_P\colon \cO_{X,f(P)}\to \cO_{Z,P}$ 是同构。
  \item $(f,f^\sharp)$ 是浸入，当且仅当 $f$ 诱导 $Z$ 与 $X$ 的某个局部闭子集的同胚，且对每个 $P\in Z$，$f^\sharp_P\colon \cO_{X,f(P)}\to \cO_{Z,P}$ 是满射。
  \item 浸入是概形范畴中的单态射。此外，浸入（resp.\ 开浸入，resp.\ 闭浸入）的复合仍是浸入（resp.\ 开浸入，resp.\ 闭浸入）。
\end{enumerate}
\end{proposition}

\begin{proposition}\label{it:1.3.16}
设 $A$ 是环。
\begin{enumerate}[label=(\roman*),leftmargin=2.8em]
  \item 对 $A$ 的每个理想 $\fa$，由典范同态 $\phi\colon A\to A/\fa$ 诱导的态射 $\varphi\colon \Spec A/\fa\to \Spec A$ 是闭浸入。
  \item $\Spec A$ 的每个闭浸入都同构于 $\Spec A/\fa\to \Spec A$，其中 $\fa$ 是 $A$ 的某个理想。
\end{enumerate}
\end{proposition}

\begin{proof}
(i) 把 $\varphi\colon \Spec A/\fa\to \Spec A$ 在底拓扑空间上诱导到 $V(\fa)$ 上的同胚一事留给读者验证。为证 $\varphi^\sharp\colon \cO_{\Spec A}\to \varphi_*\cO_{\Spec A/\fa}$ 是满射，我们使用 \ref{it:1.1.8}(iii) 与以下事实：对每个 $\fp\in \Spec A$，有
\[
(\varphi_*\cO_{\Spec A/\fa})_\fp=
\begin{cases}
(A/\fa)_{\fp/\fa}, & \text{若 } \fp\in V(\fa),\\
0, & \text{其他}.
\end{cases}
\]

(ii) 设 $\varphi\colon Z\to \Spec A$ 是闭浸入。我们先证 $Z$ 是仿射的。把 $Z$ 视为 $\Spec A$ 的子集。对任意点 $P\in Z$，可找到 $P$ 在 $\Spec A$ 中的开邻域 $U_P$，使 $U_P\cap Z$ 是 $Z$ 的仿射开子概形。选取 $f_P\in A$，使 $D(f_P)$ 是 $P$ 在 $\Spec A$ 中的邻域且包含在 $U_P$ 中。则 $D(f_P)\cap Z$ 是 $Z$ 的仿射开子概形。事实上，若 $U_P\cap Z=\Spec A'$，且 $g$ 是 $f_P$ 在复合
\[
A=\cO_{\Spec A}(\Spec A)\to \cO_{\Spec A}(U_P)\xrightarrow{\varphi^\sharp}\cO_Z(U_P\cap Z)=A'
\]
下的像，则由 \ref{it:1.3.9}(ii) 有 $D(f_P)\cap Z=D(g)\cong \Spec A'_g$。用形如 $D(f)$ 的开子集覆盖 $\Spec A-Z$。这样得到 $\Spec A$ 的覆盖 $\{D(f_i)\}_{i\in I}$，使对每个 $i$，$D(f_i)\cap Z$ 是 $Z$ 的仿射开子概形。（其中某些可能是空集。）因 $\Spec A$ 拟紧，可设此覆盖是有限的。因 $D(f_i)$（$i\in I$）覆盖 $\Spec A$，$f_i$（$i\in I$）生成单位理想。对每个 $i$，把 $f_i$ 在同态 $\varphi^\sharp\colon A=\cO_{\Spec A}(\Spec A)\to \cO_Z(Z)$ 下的像记为 $\bar f_i$。则 $\bar f_i$（$i\in I$）生成单位理想 $\cO_Z(Z)$。此外每个 $Z_{\bar f_i}=D(f_i)\cap Z$ 是仿射的。故由 \ref{it:1.3.10}，$Z$ 是仿射的。

设 $Z=\Spec B$。则闭浸入 $Z\to \Spec A$ 由同态 $\phi\colon A\to B$ 诱导。设 $\fa=\ker \phi$。我们有概形态射的交换图表：
\begin{center}
\begin{tikzcd}
\Spec B \arrow[r, "\varphi"] \arrow[d, "\theta"'] & \Spec A \\
\Spec A/\fa \arrow[ur] &
\end{tikzcd}
\end{center}
其中竖箭头 $\theta$ 由单同态 $A/\fa\to B$ 诱导。利用 $\Spec B\to \Spec A$ 与 $\Spec A/\fa\to \Spec A$ 是闭浸入这一事实，可以证明 $\theta\colon \Spec B\to \Spec A/\fa$ 也是闭浸入。由下面的 \ref{it:1.3.17}，$\varphi\colon \Spec B\to \Spec A$ 的像是 $V(\fa)$。故 $\theta\colon \Spec B\to \Spec A/\fa$ 是底拓扑空间上的同胚。因 $A/\fa\to B$ 是单射，对每个 $f\in A$，同态 $(A/\fa)_{\bar f}\to B_{\phi(f)}$ 是单射，其中 $\bar f$ 是 $f$ 在 $A/\fa$ 中的像。故
\[
\theta^\sharp(D(f))\colon \cO_{\Spec A/\fa}(D(\bar f))\to (\theta_*\cO_{\Spec B})(D(\bar f))
\]
是单射。于是 $\theta^\sharp\colon \cO_{\Spec A/\fa}\to \theta_*\cO_{\Spec B}$ 是单射。因 $\theta$ 是闭浸入，$\theta^\sharp$ 必是同构。于是闭浸入 $Z\to \Spec A$ 与 $\Spec A/\fa\to \Spec A$ 等价。
\end{proof}

\begin{lemma}\label{it:1.3.17}
设 $\phi\colon A\to B$ 是环同态，$f\colon \Spec B\to \Spec A$ 是由 $\phi$ 诱导的态射。则有 $\ol{f(\Spec B)}=V(\ker \phi)$。
\end{lemma}

\begin{proof}
我们给出三个不同的证明，以展示不同的技巧。

\emph{第一个证明}。对 $\Spec A$ 的每个子集 $Y$，我们断言 $\ol{Y}=V(\bigcap_{\fp\in Y}\fp)$。于是有
\[
\begin{aligned}
\ol{f(\Spec B)}
&=V\Bigl(\bigcap_{\fp\in f(\Spec B)}\fp\Bigr)
=V\Bigl(\bigcap_{\fq\in \Spec B}\phi^{-1}(\fq)\Bigr)
=V\Bigl(\phi^{-1}\Bigl(\bigcap_{\fq\in \Spec B}\fq\Bigr)\Bigr)\\
&=V\bigl(\phi^{-1}(\sqrt{(0)})\bigr)=V(\sqrt{\ker\phi})=V(\ker\phi).
\end{aligned}
\]
下面证明该断言。显然有 $Y\subset V(\bigcap_{\fp\in Y}\fp)$，从而 $\ol{Y}\subset V(\bigcap_{\fp\in Y}\fp)$。设 $\ol{Y}=V(\fa)$，其中 $\fa$ 是 $A$ 的某个理想。则对每个 $\fp\in Y$ 有 $\fa\subset \fp$。故 $\fa\subset \bigcap_{\fp\in Y}\fp$，从而 $V(\bigcap_{\fp\in Y}\fp)\subset V(\fa)=\ol{Y}$。

\emph{第二个证明}。显然有 $f(\Spec B)\subset V(\ker \phi)$，从而 $\ol{f(\Spec B)}\subset V(\ker \phi)$。对每个 $\fp\in V(\ker \phi)$，由 Zorn 引理，存在 $A$ 的素理想 $\fp'$，它在包含 $\ker \phi$ 且包含于 $\fp$ 的素理想族中是极小元。我们断言 $\fp'\in f(\Spec B)$。因 $\fp\in V(\fp')=\ol{\{\fp'\}}$，这蕴含 $\fp\in \ol{f(\Spec B)}$。于是 $V(\ker \phi)\subset \ol{f(\Spec B)}$。

下面证明该断言。$\fp'/\ker \phi$ 是 $A/\ker \phi$ 的极小素理想。此外，同态 $A/\ker \phi\to B$ 是单射。用 $A/\ker \phi$ 替换 $A$，问题归结为证明以下断言：对任意单同态 $\phi\colon A\to B$ 与 $A$ 的任意极小素理想 $\fp$，存在 $B$ 的素理想 $\fq$，使 $\phi^{-1}(\fq)=\fp$。注意同态 $A_\fp\to B_\fp$ 是单射。由 \ref{it:1.2.2}(i)，在 $A_\fp$ 中有 $0\ne 1$。故在 $B_\fp$ 中有 $0\ne 1$，从而 $\Spec B_\fp$ 非空。因此可找到 $B$ 的素理想 $\fq$，它与 $\phi(A-\fp)$ 不交，即 $\phi^{-1}(\fq)\subset \fp$。由 $\fp$ 的极小性，必有 $\phi^{-1}(\fq)=\fp$。

\emph{第三个证明}。同样显然有 $\ol{f(\Spec B)}\subset V(\ker \phi)$。给定 $\fp\in V(\ker \phi)$，为证 $\fp\in \ol{f(\Spec B)}$，只需证明 $\fp$ 的每个邻域 $D(a)$ 都含有 $f(\Spec B)$ 中的元素。于是需要证明：对 $A$ 的每个包含 $\ker \phi$ 的素理想 $\fp$ 与每个 $a\in A-\fp$，存在 $B$ 的素理想 $\fq$，使 $a\notin \phi^{-1}(\fq)$。只需证明 $D(\phi(a))$ 非空，等价地，$\phi(a)$ 不幂零。事实上，若 $\phi(a)$ 幂零，则 $a$ 的某个幂 $a^n$ 属于 $\ker \phi$，从而属于 $\fp$。这与 $a\notin \fp$ 矛盾。
\end{proof}
